\subsection{QUOTIENT SPACES OF COMPACT SPACES AND LOCALLY COMPACT SPACES}

PROPOSITION 8. Let X be a compact space, R an equivalence relation on X, C the graph of R in X × X, f the canonical mapping X → X/R. Then the following conditions are equivalent:

\begin{enumerate}
\def\labelenumi{\alph{enumi})}
\item
  C is closed in \(\mathbf{X} \times \mathbf{X}\) .
\item
  R is closed.
\item
  \(f\) is proper.
\item
  X/R is Hausdorff.
\end{enumerate}

Furthermore, if these conditions are satisfied then X/R is compact.

R is closed if and only if \(f\) is closed; hence b) implies c) by reason of Theorem 1b) of no. 2. That c) implies d) is a particular case of no. 1, Proposition 5, Corollary 2. d) implies a) for any topological space \(\mathbf{X}\) ( \(\S 8\) , no. 3, Proposition 8). It remains to show that a) implies b). If \(\mathbf{F}\) is a closed subset of \(\mathbf{X}\) its saturation (with respect to \(\mathbf{R}\) ) is \(\mathrm{pr}_2(\mathbf{C} \cap (\mathbf{F} \times \mathbf{X}))\) ; by hypothesis, \(\mathbf{C} \cap (\mathbf{F} \times \mathbf{X})\) is a closed subset of the compact space \(\mathbf{X} \times \mathbf{X}\) , and is therefore compact ( \(\S 9\) , no. 3, Proposition 3); the result now follows from the continuity of \(\mathrm{pr}_2\) ( \(\S 9\) , no. 4, Theorem 2, Corollary 2).

Finally it is clear that if X/R is Hausdorff then it is compact (§ 9, no. 4, Theorem 2).

PROPOSITION 9. Let X be a locally compact space, R an equivalence relation on X, C the graph of R in X × X, f the canonical mapping X → X/R; let X' be the compact space obtained by adjoining a point at infinity ω to X, and let R' be the equivalence relation on X' whose graph is C' = C ∪ \{(ω, ω)\}. Then the following conditions are equivalent:

\begin{enumerate}
\def\labelenumi{\alph{enumi})}
\item
  \(f\) is proper.
\item
  The saturation of each compact subset of X with respect to R is compact.
\item
  \(\mathbf{R}'\) is closed.
\item
  The restriction of \(\mathbf{pr}_2\) to \(\mathbf{C}\) is proper.
\item
  R is closed and the equivalence classes with respect to R are compact.
\end{enumerate}

Furthermore, if these conditions are satisfied, X/R is locally compact.

\begin{enumerate}
\def\labelenumi{\alph{enumi})}
\item
  \(\Longrightarrow\) b): Since \(\mathbf{X} / \mathbf{R} = f(\mathbf{X})\) and \(f\) is proper, \(\mathbf{X} / \mathbf{R}\) is Hausdorff (no. 1, Proposition 5, Corollary 2); hence the image under \(f\) of every compact subset \(\mathbf{K}\) of \(\mathbf{X}\) is compact (§ 9, no. 4, Theorem 2, Corollary 1). The saturation of \(\mathbf{K}\) with respect to \(\mathbf{R}\) is \(\overline{f}^1(f(\mathbf{K}))\) and is therefore compact by no. 2, Proposition 6.
\item
  \(\Longrightarrow\) c): If \(F'\) is closed in \(X'\) and does not contain \(\omega\) , then \(F'\) is a compact subset of \(X\) ; hence its saturation with respect to \(R'\) , which is the same as its saturation with respect to \(R\) , is compact and a fortiori closed in \(X'\) . If \(\omega \in F'\) and if \(F = F' \cap X = F' - \{\omega\}\) , then the saturation of \(F'\) with respect to \(R'\) is the union of \(\{\omega\}\) and the saturation \(H\) of \(F\) with respect to \(R\) ; hence it is enough to show that \(H\) is closed in \(X\) (i.e.~that \(R\) is a closed relation). For this, it is enough to show that if \(K\) is any compact subset of \(X\) then \(H \cap K\) is compact ( \(\S 9\) , no. 7, Proposition 11). Now the saturation \(L\) of \(K\) with respect to \(R\) is compact by hypothesis, and \(H \cap L\) is the saturation of \(F \cap L\) , which is also compact; a fortiori \(H \cap K = (H \cap L) \cap K\) is compact.
\item
  \(\Longrightarrow\) d): Since \(\mathbf{X}'\) is regular (§ 9, no. 2, Corollary to Proposition 1), \(\mathbf{C}'\) is closed in \(\mathbf{X}' \times \mathbf{X}'\) (§ 8, no. 6, Proposition 14) and therefore compact. It follows that \(\mathbf{C}'\) is the one-point compactification of \(\mathbf{C}\) (§ 9, no. 8, Theorem 4). Since the restriction to \(\mathbf{C}'\) of \(\mathrm{pr}_2: \mathbf{X}' \times \mathbf{X}' \to \mathbf{X}'\) is continuous at \(\omega\) , the result follows from no. 3, Corollary to Proposition 7. d) \(\Longrightarrow\) e): If \(\mathbf{F}\) is any closed subset of \(\mathbf{X}\) , then \(\mathbf{C} \cap (\mathbf{F} \times \mathbf{X})\) is closed in \(\mathbf{C}\) , whence the saturation of \(\mathbf{F}\) with respect to \(\mathbf{R}\) , which is equal to \(\mathrm{pr}_2(\mathbf{C} \cap (\mathbf{F} \times \mathbf{X}))\) , is closed in \(\mathbf{X}\) (no. 1, Proposition 1). Also the equivalence class of \(x \in \mathbf{X}\) mod \(\mathbf{R}\) is homeomorphic to the inverse image of \(\{x\}\) under the restriction of \(\mathrm{pr}_2\) to \(\mathbf{C}\) and is therefore compact {[}no. 2, Theorem 1 b){]}.
\item
  \(\Rightarrow\) a): If R is closed, then by definition \(f\) is closed, and for each \(z \in \mathbf{X} / \mathbf{R}\) , \(\overline{f}^1(z)\) is an equivalence class mod R and is therefore compact; hence \(f\) is proper by Theorem 1 b) of no. 2.
\end{enumerate}

Finally we have to prove that X/R is locally compact. \(X^{\prime}/R^{\prime}\) is compact by c) and Proposition 8; the relation R is that induced on X by \(R^{\prime}\) ; X is open in \(X^{\prime}\) and is saturated with respect to \(R^{\prime}\) ; hence X/R is homeomorphic to the image \(f^{\prime}(X)\) of X under the canonical mapping \(f^{\prime}: X^{\prime} \to X^{\prime}/R^{\prime}\) (§ 3, no. 6, Proposition 10, Corollary 1). Now \(f^{\prime}(X)\) is open in \(X^{\prime}/R^{\prime}\) , and hence is a locally compact subspace of \(X^{\prime}/R^{\prime}\) .

Q.E.D.

COROLLARY. Let X be a Hausdorff space, Y a topological space, \(f: X \rightarrow Y\) a proper mapping. Then for X to be compact (resp. locally compact) it is necessary and sufficient that \(f(X)\) is compact (resp. locally compact), and it is sufficient that Y is compact (resp. locally compact).

If X is compact (resp. locally compact) the fact that \(f(\mathbf{X})\) is compact (resp. locally compact) is a consequence of no. 1, Proposition 5, Corollary 4 and of Propositions 8 and 9 (the case where X is compact is also a consequence of no. 1, Proposition 5, Corollary 2 and § 9, no. 4, Theorem 2). Conversely, if \(Z = f(\mathbf{X})\) is compact (resp. locally compact) then since

\(f_{\mathbf{Z}}: \mathbf{X} \to f(\mathbf{X})\) is proper {[}no. 1, Proposition 3 a){]} it follows that \(\mathbf{X}\) is compact (resp. locally compact) by reason of Proposition 6 of no. 2 and Proposition 7 of no. 3. Finally, if \(\mathbf{Y}\) is compact (resp. locally compact) then so is \(f(\mathbf{X})\) which is closed in \(\mathbf{Y}\) (no. 1, Proposition 1 and § 9, no. 7, Proposition 13).

\pandocbounded{\includegraphics[keepaspectratio]{images/part001_1db83a62.jpg}}

Remark. If X is locally compact but not compact, then a closed equivalence relation R on X need not be Hausdorff (Chapter IX, § 4, Exercise 14); and even if R is Hausdorff, X/R need not be locally compact (Exercise 17). However, we have the following criterion:

PROPOSITION 10. Let X be a locally compact space, R an open Hausdorff equivalence relation on X, and let \(f: \mathrm{X} \to \mathrm{X}/\mathrm{R}\) be the canonical mapping. Then \(\mathrm{X}/\mathrm{R}\) is locally compact, and if \(\mathbf{K}'\) is any compact subset of \(\mathrm{X}/\mathrm{R}\) there is a compact subset K of X such that \(f(\mathbf{K}) = \mathbf{K}'\) .

The first assertion is a consequence of the facts that each \(x \in \mathbf{X}\) has a compact neighbourhood \(\mathbf{V}\) and that \(f(\mathbf{V})\) is a compact neighbourhood of \(f(x)\) ( \(\S 5, \text{no. } 3, \text{ Proposition } 5\) and \(\S 9, \text{ no. } 4, \text{ Theorem } 2, \text{ Corollary } 1\) ). For each \(y \in \mathbf{K}'\) let \(\mathbf{V}(y)\) be a compact neighbourhood of some point of \(\overline{f}^1(y)\) in \(\mathbf{X}\) , so that \(f(\mathbf{V}(y))\) is a compact neighbourhood of \(y\) . There are a finite number of points \(y_i \in \mathbf{K}'\) such that the \(f(\mathbf{V}(y_i))\) cover \(\mathbf{K}'\) . Let \(\mathbf{K}_1\) be the compact set \(\bigcup_i \mathbf{V}(y_i)\) in \(\mathbf{X}\) ; we have \(\mathbf{K}' \subset f(\mathbf{K}_1)\) ; hence \(\mathbf{K} = \mathbf{K}_1 \cap \overline{f}^1(\mathbf{K}')\) is compact (because it is closed in \(\mathbf{K}_1\) ) and \(f(\mathbf{K}) = \mathbf{K}'\) .

